\section{실험 및 결과}
\label{sec:results}

\paragraph{평가 프로토콜.}
하나의 경기는 여러 팀과 여러 분석 페이즈를 만들기 때문에, 팀-페이즈
행을 무작위로 나누면 같은 경기의 공간 조건과 교전 맥락이 훈련과 테스트에
동시에 들어갈 수 있다. 이를 피하기 위해 seed 42에서 경기 단위로
17,460개 훈련, 2,182개 검증, 2,182개 테스트 경기를 분할하였다.
주 비교는 테스트 경기에서 모든 모델이 공통으로 예측을 낸 134,245개의
위험 집합 셀에서 수행하였다. 여기서 하나의 셀은 \(b_k\)까지 생존했고
다음 페이즈가 관측되는 한 팀의 \((match, team, phase)\) 관측치이다.
유의성 검정은 셀 단위가 아니라 경기 단위의 성능을 짝지어 수행하여,
팀 수가 많은 경기나 긴 경기 하나가 검정을 지배하지 않도록 했다.
평가지표는 다음과 같이 정의한다. AUC는 ROC 곡선 아래 면적으로, 실제
탈락 셀에 더 높은 위험을 부여하는 순위 성능을 요약하며 높을수록 좋다.
AP는 precision-recall 곡선 아래 평균 정밀도로, 양성 셀이 상대적으로 적은
hazard 예측에서 순위 성능을 보완하며 높을수록 좋다. BCE는 실제
hazard 목표값과 예측 확률 사이의 binary cross-entropy 로그 손실이다.
Brier 점수는 예측 확률과 실제 목표값 사이의 평균 제곱오차이다. ECE는
예측 확률을 10개 동일 폭 구간으로 나누고, 각 구간의 평균 예측 위험과
실제 탈락 비율의 차이를 셀 수로 가중평균한 calibration error이다. BCE,
Brier 점수, ECE는 낮을수록 좋다.

\paragraph{비교 모델.}
특징만 사용하는 histogram gradient boosting(HGB) 기준 모델, 모든
인접 행렬을 0으로 두되 시간 구조는 유지한 그래프 없음
DynamicSTGNN, 그리고 세 그래프 변형을 비교하였다. 그래프 변형은
근접만 사용한 모델, 근접과 교전을 사용한 모델, 근접, 교전, 공동
압력을 모두 사용한 전체 그래프 모델이다. HGB는 강한 표 기반 기준 모델
역할을 하며, 그래프 없음 DynamicSTGNN은 신경망 구조와 시간 처리 방식은
유지한 채 그래프 관계만 제거한 통제 모델이다.

\begin{table}[htbp]
\centering
\caption{134,245개 정렬 테스트 셀에서의 hazard 예측 결과.}
\label{tab:main}
\scriptsize
\setlength{\tabcolsep}{3pt}
\begin{tabular}{lrrrrr}
\toprule
모델 & AUC & AP & BCE & Brier & ECE \\
\midrule
HGB 특징 & 0.7651 & 0.6626 & 0.5140 & 0.1708 & 0.0075 \\
그래프 없음 & 0.7681 & 0.6695 & 0.5090 & 0.1692 & 0.0039 \\
근접 & 0.7741 & 0.6762 & 0.5046 & 0.1674 & 0.0023 \\
근접+교전 & 0.7755 & 0.6779 & 0.5036 & 0.1670 & 0.0028 \\
전체 그래프 & \textbf{0.7782} & \textbf{0.6808} & \textbf{0.5016} & \textbf{0.1662} & \textbf{0.0015} \\
\bottomrule
\end{tabular}
\end{table}

표~\ref{tab:main}에서는 두 패턴이 나타난다. 첫째, 팀-페이즈 특징만으로도
탈락 위험은 충분히 학습 가능하다. HGB는 AUC 0.765에 도달하고, 그래프 없음
신경망 통제 모델은 0.768에 도달한다. 둘째, 그래프 관계는 작지만 일관된
이득을 더한다. 전체 그래프는 그래프 없음 모델 대비 +0.010 AUC, +0.011 AP를
개선하고, HGB 대비 +0.013 AUC, +0.018 AP를 개선한다. 경기 단위로 짝지은
Wilcoxon 검정에서 전체 그래프와 그래프 없음 모델의 차이는 유의하며(AUC
\(p=1.17\times10^{-30}\), AP \(p=3.12\times10^{-23}\)), 전체 그래프와 HGB의
차이도 유의하다(AUC \(p=2.60\times10^{-37}\), AP
\(p=9.28\times10^{-36}\)). 전체 그래프 모델은 Brier 점수 0.1662와 ECE
0.0015로 확률 품질도 가장 좋았고, 평균 예측 위험은 경험적 양성
비율에 가깝다. 주요 squad 맵별 보조 실험에서도 Baltic, Desert, Neon,
Savage, Tiger 모두 전체 그래프가 그래프 없음 모델보다 높은 AUC를 보였다
(\(\Delta\)AUC \(+0.0059\)--\(+0.0119\), AP \(+0.003\)--\(+0.020\)).

\begin{table}[htbp]
\centering
\caption{동일한 HGB 절차에서의 표현 ablation. STDB 시공간 특징은
궤적과 존/적 노출 특징을 포함한다.}
\label{tab:repr}
\small
\begin{tabular}{lrrr}
\toprule
특징 집합 & \(d\) & AUC & AP \\
\midrule
전략 점수 요약 & 5 & 0.6490 & 0.4547 \\
전략 점수 원입력 & 12 & 0.7184 & 0.5703 \\
STDB 시공간만 & 16 & 0.7124 & 0.5895 \\
기존 집계 & 57 & 0.7522 & 0.6476 \\
기존+STDB 시공간 & 72 & \textbf{0.7652} & \textbf{0.6626} \\
\bottomrule
\end{tabular}
\vspace{2pt}
\begin{minipage}{0.96\linewidth}
\scriptsize
\textit{주.} 전략 점수 원입력은 페이즈 인덱스, 피해/처치/기절, 피해를 준
상대 수, 이동 거리, 평균 속도, 100m 적 밀도, 아이템 및 보급 획득을 포함한
12개 변수이다. 전략 점수 요약은 이 원입력에서 만든 fight score, active
score, strategy type, type survival rate와 페이즈 인덱스로 구성한 5개
변수이다.
\end{minipage}
\end{table}

표~\ref{tab:repr}에서는 모델 종류가 아니라 표현 자체의 효과를 분리한다.
학습기를 전략 점수 요약 5개 특징으로 제한하면 AUC는 0.649로
떨어진다. 점수 원입력, 기존 텔레메트리 집계, STDB 이동/노출 특징을
추가하면 AUC는 0.765까지 상승한다. 즉 궤적과 존/적 노출을 포함한
시공간 특징은 기존 집계 위에서도 추가 이득을 제공하며, 가장 큰 이득은
전략 점수 요약 대비 텔레메트리 기반 페이즈 표현 자체에서 나온다
(표의 pooled AUC 기준 0.649에서 0.765로 상승).
동적 그래프는 그 위에 추가적인 관계 신호를 제공한다. 그래프 셔플 통제 실험은
이 이득이 단순한 매개변수 증가만으로 설명되기 어렵다는 보조 근거를 제공한다.
추론 시 관측 그래프 내부의 팀 identity를 섞으면 전체 그래프 모델은 5회 반복
평균 AUC 0.701로 하락하여 그래프 없음 통제 모델보다 낮아진다. 다만 어떤 간선
속성이 핵심인지를 분리하려면 차수 보존 무작위화나 관계별 순열 통제가
추가로 필요하다.
